\chapter[State of the Art]{State of the Art and Technological Background}
\begin{spacing}{1.2}
\minitoc
\thispagestyle{MyStyle}
\end{spacing}
\newpage

\section*{INTRODUCTION}
\addcontentsline{toc}{section}{INTRODUCTION}
The project presented in this report sits at the intersection of three technological landscapes: industrial quotation and Request for Quotation (RFQ) management; microservices and layered software architecture; and conversational systems built on large language models (LLMs). This chapter reviews each landscape with the aim of isolating the gap the proposed platform addresses, rather than surveying any one field exhaustively. Section~2.6 synthesises the result and states the positioning of the proposed platform.

\section{Industrial RFQ Systems and the Lifecycle Visibility Gap}
Three categories of commercial software support industrial quotation. \textbf{ERP-embedded quoting modules}, exemplified by SAP CPQ~\cite{sap_cpq_product}, integrate quote generation with the transactional and master-data backbone of the enterprise. \textbf{Stand-alone CPQ tools}, exemplified by Salesforce CPQ~\cite{salesforce_cpq_docs}, focus on configuration, pricing, and proposal generation with deeper support for configuration logic and discounting rules. \textbf{Bespoke spreadsheet-based estimation} continues to dominate in domains where the cost model embeds uncodified expert judgement --- the GHI estimation workbook of Chapter~1 is one such instance. Business process management literature~\cite{dumas2018bpm} frames the RFQ activity itself as a coordinated orchestration spanning multiple roles and decision points, distinct from any individual price calculation.

Each category optimises a different slice of the quotation activity. None is designed around the specific gap identified during the audit at \clientCompany: cross-departmental lifecycle coordination, leadership-level visibility, historical learning, early surfacing of feasibility risk, and structured capture of decisions accumulated along the way. The bottleneck is not the workbook itself but the lifecycle that surrounds it; the combination of capabilities required to address that lifecycle --- RFQ orchestration, document intelligence, evidence-grounded conversational access, and a foundation for future learning --- is what the proposed platform positions itself to deliver.

\section{Microservices Architecture and Project-Specific Layering Synthesis}
At platform level, the proposed system is structured as five responsibility-oriented microservices~\cite{newman2021microservices}, decomposed along the natural seams of the RFQ lifecycle problem: operational lifecycle management, document and workbook intelligence, conversational assistance, user interaction, and identity and access management. The detailed decomposition is developed in Chapter~4 (Section~4.2). Inside each backend service, the implementation follows the \textbf{BACAB pattern} --- a project-specific synthesis of the traditional three-tier architecture and the Ports-and-Adapters discipline~\cite{cockburn2005hexagonal} --- which separates entry points (routes), business orchestration (controllers), persistence access (datasources), external communication (connectors), and data representation (translators), with models, configuration, and utilities as support layers. The pattern preserves lifecycle traceability, separates operational truth from derived intelligence, and supports independent service evolution; it is developed in Chapter~4 (Section~4.2) and implemented in Chapter~5 (Section~5.1).

\section{Conversational AI Patterns for Evidence-Grounded Systems}
The conversational layer of the proposed platform builds on the recent generation of large language models~\cite{brown2020gpt3}. The four design patterns surveyed below share this substrate; they differ in the discipline they impose around the model's behaviour, not in the model itself. The question this section sets up is which pattern fits a cost-of-error-asymmetric industrial context.

\textbf{Naive retrieval-augmented generation (RAG)}~\cite{lewis2020rag} appends retrieved passages to the prompt; the pattern suits question-answering over static documents but offers no structural defence against the model going beyond what was retrieved. \textbf{LLM agents with autonomous tool calling}, exemplified by ReAct~\cite{yao2023react}, grant the model autonomy over tool selection and invocation order --- appropriate for exploratory contexts, but a structural concern when every retrieved value may end up in a contractual quotation. \textbf{Constrained tool-calling chatbots} built on provider function-calling APIs~\cite{anthropic_tool_use_docs} retain execution control in application code but still permit the model to influence policy through its choice of intent classification. \textbf{Trust-bound or policy-as-configuration systems}, the emerging fourth family, allocate every safety-relevant decision --- which tool may run, which user may see what, when to refuse, what evidence enters the prompt --- to deterministic platform components rather than to the language model. The field has not yet settled on a canonical reference for this family, but it is the architectural family identified by the proposed platform, for the safety and grounding reasons developed in Section~2.4 and operationalised in Chapter~4 (Section~4.4).

\section{Grounding, Hallucination, and Safety Concerns}
Any deployment of a language-model-based conversational system must contend with \textbf{hallucination}: factually incorrect content that is locally coherent and stylistically indistinguishable from correct content. The survey by Ji et al.~\cite{ji2023hallucination} establishes hallucination as a structural property of the generative regime rather than a transient quirk --- mitigation must be designed into the surrounding system rather than expected to emerge from the model alone. A related concern is \textbf{context dilution}, in which relevant evidence is buried among unrelated material in a long prompt; this concern motivates the per-path field whitelist developed in Chapter~4 (Section~4.5).

Retrieval-augmented generation reduces but does not eliminate hallucination --- the model may still drift beyond the retrieved evidence, and a system that relies on RAG alone offers no structural defence against that drift. The proposed copilot therefore combines four complementary mechanisms: \textbf{deterministic guardrails} inspect the draft before delivery (evidence grounding, scope adherence, shape conformance); an \textbf{LLM-as-Judge layer}~\cite{zheng2023judge} performs final language-level verification (fabrication, scope drift, forbidden inference) that deterministic rules cannot reliably catch; \textbf{structured outputs}~\cite{anthropic_structured_outputs_docs} constrain the model to schema-conformant JSON; and \textbf{policy-as-configuration}~\cite{bai2022constitutional} expresses safety policy as data rather than as code, enabling auditable scope expansion without pipeline changes. The combination is the architectural commitment behind the trust-boundary architecture, developed in Chapter~4 (Section~4.4) and validated in Chapter~6 (Section~6.4).

\section{Deterministic Document Intelligence for Engineering Artifacts}
The estimation workbook and Material Requisition package introduced in Chapter~1 are structured documents whose content drives the cost model the workbook produces. Two broad approaches exist for extracting information from such documents. \textbf{Deterministic parsing} exploits document structure directly: a spreadsheet with a known sheet layout is parsed cell by cell, a document with a known section structure is parsed by section identifier, and the values extracted are returned as structured data without any generative interpretation. Mature parser ecosystems exist for the principal document formats. \textbf{LLM-based document question answering} presents the document content to a language model and asks for specific values or answers; the approach is more flexible in the face of document variability but is subject to the hallucination concerns of Section~2.4, particularly in the extraction of numeric values whose correctness cannot be visually verified by the user.

The choice between these approaches is constrained by the cost-of-error posture (Chapter~1 §1.5): a wrong number propagated into a contractual quotation costs far more than a missing one that prompts a manual review. The intelligence service therefore adopts a \textbf{deterministic-parsing-first stance}: the parsers extract structured profiles, and values they cannot extract reliably are surfaced as anomaly flags rather than approximated. This design choice is the most important contribution of the document-intelligence layer, and it is demonstrated by the parser implementations in Chapter~5 (Section~5.4).

The workbook and Material Requisition package are governed by industrial standards that constrain their content. Pressure vessel construction is governed by the ASME Boiler and Pressure Vessel Code, Section VIII Division 1~\cite{asme_bpvc_viii_div1_2025}; shell-and-tube heat exchangers are governed by the standards of the Tubular Exchanger Manufacturers Association~\cite{tema_standards_11ed}. The intelligence service's parsers are designed to surface the standards references present in a Material Requisition package, providing the foundation on which standards-driven cost implications --- a pain point the audit identified as historically late to recognise (PP-12) --- can be surfaced earlier in the workflow.

\section{Positioning of the Proposed Platform}
The preceding sections have surveyed the categories of system against which the proposed platform must be positioned: ERP-embedded quoting modules, stand-alone CPQ tools, bespoke spreadsheet-based estimation, the four families of LLM-powered conversational patterns, and the document-intelligence traditions. These categories are not in direct competition with one another; each addresses a different part of the activity that surrounds an industrial quotation. The contribution of the proposed platform lies in the combination of capabilities it brings to bear on the specific gap identified during the audit, not in any claim of superiority on dimensions where the surveyed categories were never designed to compete.

The proposed platform is positioned as an \textbf{RFQ Lifecycle Intelligence Platform with a trust-bound conversational layer}. It combines lifecycle awareness inherited from the workflow and BPM tradition (Section~2.1)~\cite{dumas2018bpm}; multi-actor coordination expressed in a domain model derived from Domain-Driven Design~\cite{evans2003ddd} and developed in Chapter~4 (Section~4.3); evidence grounding inspired by safety concerns~\cite{ji2023hallucination, bai2022constitutional, zheng2023judge}; scope-bounded refusal supported by structured-output discipline~\cite{anthropic_structured_outputs_docs} and by the deterministic-guardrail family; policy-as-configuration auditability characteristic of the trust-bound architectural family; and a deterministic-parsing-first stance over engineering documents aligned with the cost-of-error posture. Chapter~4 develops the architecture that delivers this combination; Chapter~5 presents its implementation; Chapter~6 validates the architectural invariants under measured behaviour.

\section*{CONCLUSION}
\addcontentsline{toc}{section}{CONCLUSION}
This chapter has reviewed the three landscapes against which the proposed platform must be positioned and isolated the gap each surveyed category leaves unfilled relative to the audit findings of Chapter~1. The platform is positioned in that gap as an RFQ Lifecycle Intelligence Platform with a trust-bound conversational layer; it is not positioned as an ERP replacement, a CPQ competitor, a generic chatbot, or a BOQ automator. The architecture that delivers this positioning is the subject of the next chapter.